% Edition 0.5: prospective confinement, precise carrier distinctions and proof status.
\chapter[Conditional Oscillator Exclusion and Spectrum]{Conditional Oscillator Exclusion, High Spin Type, and the Prospective Physical Spectrum}
\label{ch:spectral}

\section{The same obstruction seen from symmetry and analysis}
Chapter~\ref{ch:spin-geometry} isolated a concrete ninth-axis
stabilizer of the source spatial Spin(9) group. Chapters~\ref{ch:energy}
and~\ref{ch:compactness} isolated the bound on the first radial moment
that would suffice to control global $L^2$ tails. This chapter
investigates a proposed join between these stories. The potential
vanishes along rank-one configurations, but transverse oscillations
around a rank-one valley have an energy proportional to the radial
scale. A sufficiently high spatial representation type might prohibit
the lowest transverse oscillator levels, leaving a positive linear
energy that defeats the fermionic negative contribution.

There are three distinct questions: (a) is the representation-theoretic
exclusion valid in the \emph{chosen} Hilbert carrier? (b) can the
inequality be proved for the \emph{original physical} smooth core and
closed form, rather than an auxiliary model? (c) are the relevant
physical isotypic sectors nonzero and infinite-dimensional? The first
has a substantial conditional analytical solution in an
operator-valued carrier; the second and third remain source-binding
obligations. They must not be collapsed into a single spectral claim.

\section{Transverse oscillators in nine spatial directions}
Choose a nonzero one-color row $u\in E=\mathbb R^9$ with $r=\|u\|$
and unit direction $a=u/r$. Decompose the other two color rows
\begin{equation}\label{eq:spectral-onecolor}
  x^2=t_2a+z_2,\qquad x^3=t_3a+z_3,
  \qquad z_2,z_3\in a^\perp\cong\mathbb R^8.
\end{equation}
The transverse variables comprise sixteen real coordinates. A source
of the required order-of-growth lower bound is the ordinary harmonic
oscillator
\begin{equation}\label{eq:transverse-oscillator}
  h_{\beta,r}:=-\beta\Delta_{z}+r^2\|z\|^2,
  \qquad \beta>0.
\end{equation}
In the conventional self-adjoint oscillator realization, each
coordinate contributes $\sqrt\beta r(2n_j+1)$, so the eigenvalue at
total excitation number $n=\sum_jn_j$ is
\begin{equation}\label{eq:oscillator-level}
  E_n(r)=2\sqrt\beta\,r(n+8).
\end{equation}
This is a standard theorem about the auxiliary transverse oscillator,
not an identification with the physical BFSS Hamiltonian. In
particular, the projection to a low transverse energy band is not
automatically a reducing projection for the full interacting form.

\section{The stabilizer and the branching inequality}
Fix the original spatial unit vector $a$. The source subgroup
$K_a\subset\operatorname{Spin}(9)$ is its stabilizer. The accepted
covering/coordinate theorems identify its source two-sign covering of
the corresponding orthogonal fixed-axis subgroup, and identify the
latter with $SO(8)$. For a branching argument one needs more:
a vector-compatible, correctly normalized identification of the
actual source torus and stabilizer representations with the standard
$B_4\to D_4$ branching data. Source geometry alone does not perform
this representation identification.

Under that additional hypothesis, if a dominant Spin(9) highest
weight is $\lambda=(\lambda_1,\lambda_2,\lambda_3,\lambda_4)$
and a Spin(8) constituent has highest weight $\nu$, the standard
interlacing rule entails
\begin{equation}\label{eq:branching-inequality}
  \nu_1\ge\lambda_2.
\end{equation}
The four coordinates must be treated with their exact integral or
half-integral weight lattice; the center and the two half-spinor
representations cannot be identified away by passing to $SO(8)$.

\section{The six-versus-twelve weight trap}
A particularly valuable negative discovery was that the first
fermionic weight bound depends on whether the fiber contains
\emph{Fock vectors} or \emph{Fock endomorphisms}.

On the original 48-real-label module, a spatial plane rotation acts
in twenty-four real rotation planes with half-angle frequency.
For a suitable CAR implementer of the Lie generator, each plane
contributes a fermionic Hermitian weight $\pm\frac14$.
Thus the first-weight interval on the Fock-vector module is
$[-6,6]$ in the contemplated canonical lift.
But an endomorphism fiber is $\operatorname{End}_{\mathbb C}\Fermion
\simeq\Fermion\otimes\Fermion^*$, so the action by conjugation
contains \emph{both} output and dual-input weights. Its first weight
can therefore reach $12$, not just $6$.

\begin{proposition}[Conditional low-level exclusion; operator-valued
fiber]\label{prop:operator-lowlevel}
Assume a continuous, source-compatible Spin(9) action on
$L^2(E^3;\operatorname{End}_{\mathbb C}\Fermion)$, the correct
Spin(8) stabilizer and torus identification, and the standard
branching rule \eqref{eq:branching-inequality}. Suppose a smooth
compactly supported vector belongs to a genuine Spin(9) isotypic
component of highest weight $\lambda$. Then for $m\ge0$ the strict
condition
\begin{equation}\label{eq:op-threshold}
   \lambda_2>m+12
\end{equation}
excludes transverse oscillator levels $0,\ldots,m$ on each smooth
slice of the isotypic component.
\end{proposition}
\begin{proof}[Conditional analytical argument]
The excitation-$n$ transverse oscillator carries Spin(8) first
highest weight at most $n$; tensoring it with the endomorphism
fiber adds at most $12$, yielding $\nu_1\le n+12$. But any Spin(8)
constituent of the specified Spin(9) isotype must satisfy
$\nu_1\ge\lambda_2$. If $\lambda_2>m+12$, these inequalities
are incompatible for every $n\le m$.

The pointwise restriction of an arbitrary $L^2$ equivalence class to
a single slice is generally meaningless. The derivation therefore
works on a smooth compactly supported \emph{isotypic} vector, where
the finite-dimensional group orbit span makes the stabilizer
intertwiner precise; only afterwards can a separate form-core
argument attempt a closed-domain extension.
\end{proof}
\formalrecord{\path{contributions/bfss-kotler/BFSS_SPECTRAL_WHOLE_SECTOR_OSCILLATOR_EXCLUSION_2026_11_07.md}.
This is a \textbf{conditional analytical theorem}, not a strict Lean
acceptance on the exact original physical Fock-vector space. Atlas C-01.}

For the \emph{original vector} fiber, the proposed corresponding
sufficient threshold would be $\lambda_2>m+6$, but it requires an
actual global Fock-vector lift and its compatibility with the
physical gauge action and closed charge operators. That lift cannot
be inferred merely from the already accepted star automorphism of
the \emph{operator} algebra. Nor may the operator-valued threshold
$m+12$ be replaced by $m+6$ for numerical convenience.

\section{The proposed global coercivity mechanism}
Suppose one obtains, for the \emph{actual} physical smooth core in a
specified high-type sector, a transverse oscillator lower bound
combined with a source-typed fermionic estimate
\eqref{eq:B-linear-target}. A natural target inequality is
\begin{equation}\label{eq:sector-coercivity-target}
 q_{\mathrm{phys}}(\psi)+b\|\psi\|_2^2
 \ge a\|\nabla\psi\|_2^2+
 c\int_{\mathbb R^{27}}\!\|X\|\,\|\psi(X)\|^2\,dX,
 \qquad a,c>0.
\end{equation}
All constants, projections and domain statements must refer to
\emph{that same} closed physical form. This is currently a proof
target, not a result already established by the generic energy
identity or abstract compactness theorem.

The November 6 A/B/E audits identify a promising conservative
calculation. One devotes $0<\beta<1/2$ of the kinetic coefficient
to the one-color transverse oscillator. Averaging over the actual
color sphere and using source-compatible spatial rotations gives,
\emph{conditionally}, a positive radial contribution with coefficient
\begin{equation}\label{eq:sector-margin}
  \sqrt\beta\,(m+9)-C_B,
\end{equation}
provided oscillator levels up to $m$ are genuinely absent and the
original fermionic term is bounded by $C_B\|X\|$. For the proposed
original source finite-CAR constant $C_B=8\sqrt3$, the illustrative
choice $\beta=0.49$ and $m=11$ gives the strictly positive margin
\begin{equation}\label{eq:sector-example-margin}
 \sqrt{0.49}(11+9)-8\sqrt3=14-8\sqrt3>0.
\end{equation}
The positivity follows from $14^2=196>192=(8\sqrt3)^2$.
The \emph{physical Fock-vector} threshold would be
$\lambda_2>17$ if the weight-six source-lift hypothesis holds;
the \emph{operator-valued} hypothesis would instead require
$\lambda_2>23$ for the same $m$. These are two different
conditional settings. Equation~\eqref{eq:sector-example-margin}
is not an accepted theorem of the original physical Hamiltonian.

\section{What a valid physical high-type sector must supply}
A mathematical projection onto a formal highest weight is not
enough. A candidate positive-eigenvalue argument needs a nonzero
closed \emph{physical} isotypic sector, invariance under the
\emph{same} Hamiltonian, and an infinite-dimensional sector to
obtain infinitely many discrete eigenvalues from compact resolvent.
The historical polynomial proposal uses gauge-invariant combinations
of spatial complex coordinates and radial bumps; its weights and
nonvanishing have been tested analytically, but a complete source-
exact, all-domain sector theorem remains separate.

\begin{openproblem}[Exact physical sector and coercivity]
Construct the original global source Spin(9) action on the physical
Fock-vector $L^2$ Hilbert space, with correct central signs and true
SU(2) gauge compatibility. Prove covariance of all sixteen literal
supercharges on the original invariant smooth core and extend it
to the closed physical charge form and associated Hamiltonian.
Then produce at least one nonzero infinite-dimensional high-type
physical sector and establish \eqref{eq:sector-coercivity-target}
there, without replacing the Hamiltonian.
\end{openproblem}

\section{The exact compactness-to-eigenpair implication}
The accepted weighted-closure theorem from
Chapter~\ref{ch:compactness} has a powerful \emph{conditional}
consequence. Assume that for a reducing sector $\mathcal H_\lambda$
we have a densely defined lower-bounded closed quadratic form
$q_\lambda$ of the \emph{original} Hamiltonian. Assume also that its
form-domain unit ball maps continuously into the existing weighted
Sobolev closure with uniform $L^2$, gradient and
$\sqrt{\|X\|}$-weighted bounds. Then Chapter~\ref{ch:compactness}
implies compact embedding of the form domain into
$\mathcal H_\lambda$.

\begin{proposition}[Conditional compact resolvent and point spectrum]
\label{prop:conditional-spectrum}
Under those assumptions, the self-adjoint operator associated with
$q_\lambda$ has compact resolvent. If the reducing sector is
infinite-dimensional, its spectrum consists of discrete eigenvalues
of finite multiplicity tending to $+\infty$, and therefore contains
positive eigenvalues. Their inclusion as eigenvalues of the original
physical Hamiltonian additionally uses the verified reducing-sector
identification. No positive gap of the full BFSS spectrum is implied.
\end{proposition}
\begin{proof}[Standard closed-form spectral argument]
A lower-bounded closed form equips its domain with a Hilbert graph
norm. The assumed compact embedding makes the form resolvent
$\mathcal H_\lambda\to\mathcal H_\lambda$ compact through the
usual variational solution map and form-domain injection.
The spectral theorem for compact self-adjoint resolvents gives
purely discrete finite-multiplicity spectrum with eigenvalues
tending upward in an infinite-dimensional sector. A reducing
subspace embeds the sector eigenvectors into the original Hilbert
operator domain. None of these general implications proves the
hypotheses on the original source.
\end{proof}

\section{Hostile escape states as scientific tests}
The most valuable prospective counterexample would be a sequence
$\psi_j$ in one fixed nonzero physical Spin(9) sector with
$\|\psi_j\|_2=1$, support escaping to radial infinity, and
$q_{\mathrm{phys}}(\psi_j)$ uniformly bounded. Such a sequence
would contradict a claimed coercive first-moment inequality in that
sector. A rank-one classical configuration alone is not such a
sequence; it does not satisfy normalization, gauge invariance,
sector membership or operator-domain conditions.

The proposed near-valley localization route is higher-risk because
cutoffs may destroy the required isotypic or gauge properties and
introduce non-negligible IMS derivative terms. We therefore prioritize
the global color-direction averaging route, but preserve any valid
adversarial escape construction instead of modifying the sector or
Hamiltonian after the fact.

\formalrecord{Source controls: \path{BFSS_SPECTRAL_CONFINEMENT_AB_AUDIT_2026_11_06.md},
\path{BFSS_SPECTRAL_CONFINEMENT_E_GLOBAL_AVERAGING_2026_11_06.md},
\path{BFSS_SPECTRAL_CONFINEMENT_ROADMAP_2026_11_06.md}, and the
conditional \path{BFSS_SPECTRAL_WHOLE_SECTOR_OSCILLATOR_EXCLUSION_2026_11_07.md}.
Open atlas obligations O-01 through O-05 and C-01/C-02.}
